%!TEX program = xelatex
\documentclass[11pt,a4paper]{article}

% ---------- 中文支持（xelatex 原生） ----------
\usepackage{fontspec}
\usepackage{xeCJK}
\setCJKmainfont{Songti SC}[BoldFont=Heiti SC]
\setCJKsansfont{Heiti SC}
\setCJKmonofont{STFangsong}

% ---------- 页面 ----------
\usepackage[margin=2.5cm]{geometry}
\usepackage{setspace}
\onehalfspacing

% ---------- 图表 ----------
\usepackage{booktabs}
\usepackage{multirow}
\usepackage{longtable}
\usepackage{graphicx}
\usepackage{caption}
\usepackage{subcaption}

% ---------- 其他 ----------
\usepackage{hyperref}
\usepackage{xcolor}
\usepackage{enumitem}
\usepackage{float}

\hypersetup{
    colorlinks=true,
    linkcolor=blue,
    urlcolor=blue,
    citecolor=blue,
}

\title{\textbf{各国经济对人工智能依赖度的文本分析}\\
\large 基于 Wikipedia 经济词条的词向量方法}
\author{Siyao Zheng\thanks{计算社会科学；上海交通大学}}
\date{\today}

\begin{document}

\maketitle

\begin{abstract}
\noindent 本文利用 30 个国家 Wikipedia \guillemotleft Economy of X\guillemotright 词条的导言文本，
通过词向量（sentence embeddings）方法构建了\guillemotleft 经济对人工智能依赖度\guillemotright 指标。
我们使用 Sentence-BERT（all-MiniLM-L6-v2）将所有句子编码为 384 维向量，并构建了涵盖
AI 技术核心、数字经济转型、高技能劳动力和数字基础设施四个维度的 23 个锚定短语。
以锚定向量的均值作为 AI 概念向量，计算各句的余弦相似度，再按国家聚合得到均值排名。
结果显示，中国、南非、印度、韩国、印度尼西亚在 AI 依赖度指标上得分最高；
沙特阿拉伯、智利、俄罗斯、德国、瑞士得分最低。该方法为课堂演示 text-as-data
分析提供了可复现的完整工作流。
\end{abstract}

\section{引言}

\textbf{研究问题：}各国经济在多大程度上表现出对人工智能及相关数字技术的依赖？
这一问题对于理解全球技术竞争格局尤为关键。传统回答依赖专利数量、
R\&D 支出或 ICT 出口占比等结构化指标；本文尝试一种互补的文本方法：
从各国经济描述的\emph{语言特征}中推断 AI 相关度。

\section{数据}

\subsection{语料}

我们从 Wikipedia API 抓取了 30 个代表性国家的\guillemotleft Economy of X\guillemotright 词条导言
（introductory extract），覆盖北美、欧洲、东亚、南亚、东南亚、中东、非洲、
拉丁美洲等地区，以及不同收入水平。语料基本信息见表~\ref{tab:corpus}。

\begin{table}[H]
\centering
\caption{语料基本信息}
\label{tab:corpus}
\begin{tabular}{lr}
\toprule
\textbf{指标} & \textbf{数值} \\
\midrule
国家数 & 30 \\
原始文档总字符数 & 74,524 \\
原始总词数（英文） & 11,606 \\
平均词数/文档 & 387 \\
\bottomrule
\end{tabular}
\end{table}

\subsection{预处理}

预处理流程如下：
\begin{enumerate}[nosep]
    \item \textbf{分句}：NLTK \texttt{sent\_tokenize}（英语）；
    \item \textbf{分词 + 小写}：NLTK \texttt{word\_tokenize} $\to$ \texttt{.lower()}；
    \item \textbf{去停用词}：NLTK 英文停用词表 + 手动添加低信息词（also, however 等）；
    \item \textbf{词形还原}：NLTK WordNetLemmatizer；
    \item \textbf{过滤}：去除纯标点/纯数字 token、长度 $\le$ 2 的 token、token 数 $<$ 3 的句子。
\end{enumerate}

预处理后得到 \textbf{588} 条有效句子，总计 6,437 个 token，来自 1,735 个唯一词型。

\section{方法}

\subsection{AI 锚定短语与概念向量}

我们设计了 \textbf{23 个 AI 锚定短语}（见表~\ref{tab:anchors}），覆盖 AI 技术核心、
经济结构转型、劳动力与技能、数字基础设施四个维度。使用 Sentence-BERT 的
\texttt{all-MiniLM-L6-v2} 模型（384 维）将所有锚定短语编码为向量，
取均值并归一化，作为 \emph{AI 概念向量} $\mathbf{v}_{\text{AI}}$。

\begin{table}[H]
\centering
\caption{AI 锚定短语及其维度}
\label{tab:anchors}
\small
\begin{tabular}{ll}
\toprule
\textbf{维度} & \textbf{锚定短语} \\
\midrule
\multirow{6}{*}{AI 技术核心}
    & artificial intelligence and machine learning technologies \\
    & deep learning neural networks and AI research \\
    & robotics and intelligent automation systems \\
    & advanced algorithms and data-driven decision making \\
    & automated manufacturing and smart factories \\
    & AI-powered software and digital platforms \\
    & semiconductor chips and AI computing hardware \\
\midrule
\multirow{5}{*}{经济结构转型}
    & digital economy and technology-driven growth \\
    & innovation in high-tech industries \\
    & knowledge economy and research development \\
    & technology exports and intellectual property \\
    & digital transformation of traditional industries \\
\midrule
\multirow{4}{*}{劳动力与技能}
    & highly skilled technology workforce \\
    & STEM education and technical training \\
    & automation replacing routine jobs \\
    & digital skills and computer literacy \\
\midrule
\multirow{7}{*}{数字基础设施}
    & internet connectivity and digital infrastructure \\
    & cloud computing and data centers \\
    & broadband access and mobile technology \\
    & e-government and digital public services \\
    & fintech and digital financial services \\
    & startup ecosystem and venture capital in tech \\
    & investment in research and development R\&D \\
\bottomrule
\end{tabular}
\end{table}

\subsection{句子编码与相似度计算}

对 588 条句子逐一编码得到 $\mathbf{s}_i$，计算每条句子与 AI 概念向量的余弦相似度：

\begin{equation}
\text{Sim}_i = \frac{\mathbf{s}_i \cdot \mathbf{v}_{\text{AI}}}{\|\mathbf{s}_i\| \cdot \|\mathbf{v}_{\text{AI}}\|}
\end{equation}

\subsection{国家级别聚合}

对每个国家 $c$ 的所有句子相似度取均值作为 AI 依赖度指标：

\begin{equation}
\text{AI-Dep}_c = \frac{1}{N_c} \sum_{i \in c} \text{Sim}_i
\end{equation}

同时报告中位数、最大值、标准差、P75、P90 以及高相似度句子占比（$\text{Sim}_i > 0.3$）。

\section{结果}

\subsection{整体分布}

表~\ref{tab:dist} 报告了 588 条句子相似度的整体分布特征。

\begin{table}[H]
\centering
\caption{句级 AI 相似度整体分布}
\label{tab:dist}
\begin{tabular}{lr}
\toprule
\textbf{指标} & \textbf{数值} \\
\midrule
句子总数 & 588 \\
相似度均值 & 0.192 \\
相似度中位数 & 0.188 \\
相似度标准差 & 0.086 \\
最小值 & $-0.031$ \\
最大值 & 0.546 \\
高相似度句子占比（$>$0.3）& 9.35\% \\
\bottomrule
\end{tabular}
\end{table}

\subsection{国家排名}

表~\ref{tab:ranking} 报告了 30 国的 AI 依赖度排名（仅 Top-10 和 Bottom-5；
完整排名见附录~\ref{sec:full}）。

\begin{table}[H]
\centering
\caption{各国 AI 依赖度排名（摘要）}
\label{tab:ranking}
\begin{tabular}{rllrrrr}
\toprule
\textbf{排名} & \textbf{国家} & & \textbf{均值} & \textbf{中位数} & \textbf{最大值} & \textbf{标准差} \\
\midrule
1 & China & 中国 & 0.258 & 0.233 & 0.546 & 0.087 \\
2 & South Africa & 南非 & 0.247 & 0.237 & 0.408 & 0.104 \\
3 & India & 印度 & 0.237 & 0.233 & 0.429 & 0.077 \\
4 & Indonesia & 印度尼西亚 & 0.231 & 0.232 & 0.375 & 0.074 \\
5 & South Korea & 韩国 & 0.229 & 0.218 & 0.454 & 0.095 \\
6 & United States & 美国 & 0.219 & 0.213 & 0.407 & 0.081 \\
7 & United Kingdom & 英国 & 0.218 & 0.219 & 0.433 & 0.075 \\
8 & Turkey & 土耳其 & 0.213 & 0.184 & 0.460 & 0.095 \\
9 & Japan & 日本 & 0.211 & 0.215 & 0.323 & 0.053 \\
10 & Singapore & 新加坡 & 0.211 & 0.215 & 0.266 & 0.043 \\
\midrule
\multicolumn{7}{c}{\dots} \\
\midrule
26 & Switzerland & 瑞士 & 0.152 & 0.131 & 0.227 & 0.053 \\
27 & Germany & 德国 & 0.149 & 0.144 & 0.334 & 0.071 \\
28 & Russia & 俄罗斯 & 0.147 & 0.150 & 0.279 & 0.069 \\
29 & Chile & 智利 & 0.141 & 0.131 & 0.280 & 0.064 \\
30 & Saudi Arabia & 沙特阿拉伯 & 0.105 & 0.094 & 0.244 & 0.071 \\
\bottomrule
\end{tabular}
\end{table}

\section{讨论}

\subsection{主要发现}

\textbf{第一名——中国}（均值 0.258，最大 0.546）：得分最高的句子来自高科技制造业转型的描述：
``Manufacturing has been transitioning toward high-tech industries such as electric vehicles,
renewable energy, telecommunications\dots'' 这意味着中国经济描述的 AI 相关语义信号不仅广泛，
而且强度高。

\textbf{新兴经济体的突出表现}：Top-5 中的南非、印度、印度尼西亚均为发展中国家。
这可能反映其经济描述中更强调\emph{转型叙事}——从传统到数字化的语义跃迁比成熟发达经济体更具文本显著性。

\textbf{发达经济体的\guillemotleft 低调\guillemotright}：美国（第 6）、英国（第 7）、日本（第 9）
排名靠前但并非最高，德国、瑞士居后。可能原因有二：(1) 发达经济体的经济描述涵盖更广泛议题，
AI/数字话题被稀释；(2) 对于 AI/数字经济，发达经济体的文本\guillemotleft 理所当然\guillemotright
而非\guillemotleft 转型目标\guillemotright，语义上不那么凸显。

\subsection{方法局限}

\begin{enumerate}[nosep]
    \item \textbf{语料规模小}：仅 30 个文档，588 句。扩展至更多国家或全词条正文可提高信噪比。
    \item \textbf{锚点短语的主观性}：23 个短语虽覆盖四个维度，但选择存在研究者自由度；
    未来的稳健性检查可做词汇替代和敏感性检验。
    \item \textbf{依赖 Wikipedia}：Wikipedia 的编辑规范和覆盖偏差可能影响跨国可比性。
    \item \textbf{语言限制}：仅分析英文词条，非英语国家的 Wikipedia 本地语言版本可能蕴含不同信息。
    \item \textbf{聚合方式简单}：均值聚合忽略了句间权重差异（如首句通常更加总体性），
    未来可考虑 tf-idf 加权或基于文档结构的加权方案。
\end{enumerate}

\section{结论}

本文使用一种低门槛、完全可复现的词向量流程，从 Wikipedia 英文经济词条中提取了
\guillemotleft 经济对 AI 依赖度\guillemotright 的国别交叉截面指标。主要发现是：

\begin{enumerate}[nosep]
    \item 中国、南非、印度在 AI 依赖度上排名最高，其经济描述中最频繁触及 AI、
    数字经济和科技转型议题；
    \item 新兴经济体在 AI 文本信号上可能比部分发达经济体更强——这可能反映了
    \emph{叙事层面}上的转型程度而非\emph{实际层面}上的技术能力；
    \item 该方法可作为课堂 text-as-data 分析的教学演示，展示从原始文本到可解释指标
    的完整 pipeline。
\end{enumerate}

\newpage
\appendix
\section{完整 30 国排名}
\label{sec:full}

\begin{longtable}{rllrrrr}
\caption{30 国 AI 依赖度完整排名} \\
\toprule
\textbf{排名} & \textbf{国家} & & \textbf{均值} & \textbf{中位数} & \textbf{最大值} & \textbf{标准差} \\
\midrule
\endfirsthead
\multicolumn{7}{c}{（续前表）} \\
\toprule
\textbf{排名} & \textbf{国家} & & \textbf{均值} & \textbf{中位数} & \textbf{最大值} & \textbf{标准差} \\
\midrule
\endhead
\bottomrule
\endfoot
1 & China & 中国 & 0.258 & 0.233 & 0.546 & 0.087 \\
2 & South Africa & 南非 & 0.247 & 0.237 & 0.408 & 0.104 \\
3 & India & 印度 & 0.237 & 0.233 & 0.429 & 0.077 \\
4 & Indonesia & 印度尼西亚 & 0.231 & 0.232 & 0.375 & 0.074 \\
5 & South Korea & 韩国 & 0.229 & 0.218 & 0.454 & 0.095 \\
6 & United States & 美国 & 0.219 & 0.213 & 0.407 & 0.081 \\
7 & United Kingdom & 英国 & 0.218 & 0.219 & 0.433 & 0.075 \\
8 & Turkey & 土耳其 & 0.213 & 0.184 & 0.460 & 0.095 \\
9 & Japan & 日本 & 0.211 & 0.215 & 0.323 & 0.053 \\
10 & Singapore & 新加坡 & 0.211 & 0.215 & 0.266 & 0.043 \\
11 & Vietnam & 越南 & 0.210 & 0.211 & 0.311 & 0.083 \\
12 & Egypt & 埃及 & 0.203 & 0.201 & 0.280 & 0.053 \\
13 & Ethiopia & 埃塞俄比亚 & 0.198 & 0.188 & 0.369 & 0.086 \\
14 & Brazil & 巴西 & 0.194 & 0.204 & 0.457 & 0.086 \\
15 & France & 法国 & 0.191 & 0.180 & 0.333 & 0.060 \\
16 & Nigeria & 尼日利亚 & 0.189 & 0.185 & 0.326 & 0.079 \\
17 & Bangladesh & 孟加拉国 & 0.188 & 0.196 & 0.337 & 0.077 \\
18 & Thailand & 泰国 & 0.188 & 0.173 & 0.522 & 0.100 \\
19 & Italy & 意大利 & 0.187 & 0.162 & 0.355 & 0.077 \\
20 & Poland & 波兰 & 0.174 & 0.142 & 0.351 & 0.077 \\
21 & Australia & 澳大利亚 & 0.169 & 0.169 & 0.293 & 0.059 \\
22 & Mexico & 墨西哥 & 0.167 & 0.141 & 0.371 & 0.088 \\
23 & Canada & 加拿大 & 0.165 & 0.150 & 0.331 & 0.082 \\
24 & Argentina & 阿根廷 & 0.154 & 0.137 & 0.237 & 0.069 \\
25 & Netherlands & 荷兰 & 0.153 & 0.148 & 0.413 & 0.107 \\
26 & Switzerland & 瑞士 & 0.152 & 0.131 & 0.227 & 0.053 \\
27 & Germany & 德国 & 0.149 & 0.144 & 0.334 & 0.071 \\
28 & Russia & 俄罗斯 & 0.147 & 0.150 & 0.279 & 0.069 \\
29 & Chile & 智利 & 0.141 & 0.131 & 0.280 & 0.064 \\
30 & Saudi Arabia & 沙特阿拉伯 & 0.105 & 0.094 & 0.244 & 0.071 \\
\end{longtable}

\section{生成统计}

\noindent
\textbf{生成日期}：2026 年 7 月 4 日 \\
\textbf{模型}：Sentence-BERT \texttt{all-MiniLM-L6-v2}（384 维） \\
\textbf{数据处理脚本}：
\texttt{preprocess\_text\_corpus.py}, \texttt{ai\_dependency\_keyword.py}, \texttt{ai\_dependency\_embedding.py} \\
\textbf{原始数据}：\texttt{data/raw/text\_corpus/economy\_corpus.jsonl} \\
\textbf{预处理输出}：\texttt{data/processed/text\_corpus/} \\
\textbf{分析结果}：\texttt{reports/ai\_dependency\_embedding.csv}

\vspace{1cm}
\noindent\textit{本报告所有数据处理与分析均可通过 pipeline 复现，无需手动操作。}

\end{document}
